\chapter{Online Learning and Drift}\label{ch:ml:online-learning-and-drift}

A short-horizon model is retrained every weekend. On a Tuesday the relation it learned changes (a venue alters its fee
schedule, a large participant changes how it trades, a regime ends) and the model keeps trading the old relation until
the next weekend. A model that updated itself every hour would have noticed on Tuesday afternoon; a model that updated
itself every trade would have spent the whole year chasing noise. This chapter is about the choice between those two
failures: how fast a model should forget, how to detect that the world has changed, and when to retrain. On a
synthetic stream whose coefficients change at random times, a detector calibrated to one false alarm a year finds a
reversal of the signal about twenty steps after it happens, and the best memory for a recursive estimator grows from 200
observations to all of them as the regimes lengthen from a hundred steps to ten thousand.

\section{Kinds of drift}

\begin{definition}[Concept drift, covariate shift]\label{def:ml:online-learning-and-drift:drift}
\emph{Concept drift}\index{concept drift} is a change over time in the conditional distribution of the target given the
features, $\P_t(y\mid x)$: the relation the model learned stops holding. \emph{Covariate shift}\index{covariate shift}
is a change in the distribution of the features, $\P_t(x)$, with $\P(y\mid x)$ unchanged.
\end{definition}

\omterm{def:ml:online-learning-and-drift:drift}{Covariate shift} hurts a model that has learned $\P(y\mid x)$ only where training data were dense and must extrapolate
elsewhere; a correctly specified model is untouched. \omterm{def:ml:online-learning-and-drift:drift}{Concept drift} hurts every model and is what markets produce:
crowding weakens a signal (Book~7, chapter~28), a structural break reverses one, a change of market structure rewires the
relation between order flow and prices. Book~7 measured decay and breaks after the fact (CUSUM tests, the signal
half-life); this chapter adapts to them as they happen.

The chapter's stream has 20\,000 steps and five features; the target is $y_t = \beta_t^\top x_t + \varepsilon_t$ with a
signal worth 5\% of the variance, and the coefficients are redrawn at regime ends whose spacing is exponential with
mean $D$ steps. The truth's own R-squared, the best any model can do, is 4.9 to 5.4\%.

\section{Recursive estimators and forgetting}

\begin{definition}[Online learning, recursive least squares, forgetting factor]\label{def:ml:online-learning-and-drift:rls}
\emph{Online learning}\index{online learning} updates a model with each new observation, in constant time and memory,
and is scored prequentially: each observation is predicted before it is learned from. \emph{Recursive least
squares}\index{recursive least squares} (RLS) updates the weighted least-squares solution
$\hat\beta_t = \arg\min_b\sum_{s\le t}\lambda^{t-s}(y_s - x_s^\top b)^2$ by
\[
k_t = \frac{P_{t-1}x_t}{\lambda + x_t^\top P_{t-1}x_t},\quad \hat\beta_t = \hat\beta_{t-1} + k_t(y_t - x_t^\top\hat\beta_{t-1}),\quad
P_t = \frac{1}{\lambda}\bigl(P_{t-1} - k_tx_t^\top P_{t-1}\bigr),
\]
where the \emph{forgetting factor}\index{forgetting factor} $\lambda\in(0,1]$ discounts old observations; their weights sum
to about $1/(1-\lambda)$, the estimator's effective memory.
\end{definition}

\begin{proposition}[Forgetting is a random-walk Kalman filter]\label{prop:ml:online:kalman}
RLS with \omterm{def:ml:online-learning-and-drift:rls}{forgetting factor} $\lambda$ is the Kalman filter (Book~4, chapter~19) for the state-space model $y_t =
x_t^\top\beta_t + \varepsilon_t$, $\beta_t = \beta_{t-1} + w_t$, in which the prediction step inflates the state covariance
by $1/\lambda$ instead of adding the covariance of $w_t$: $P_{t\mid t-1} = P_{t-1\mid t-1}/\lambda$, with unit observation
noise (Sayed and Kailath, 1994).
\end{proposition}

\begin{proof}
With $P_{t\mid t-1} = P_{t-1}/\lambda$ and unit noise variance, the Kalman gain is $P_{t\mid t-1}x_t/(1 + x_t^\top P_{t\mid
t-1}x_t) = P_{t-1}x_t/(\lambda + x_t^\top P_{t-1}x_t) = k_t$; the state update is RLS's; the filtered covariance
$(I - k_tx_t^\top)P_{t\mid t-1} = (P_{t-1} - k_tx_t^\top P_{t-1})/\lambda$ is RLS's $P_t$.
\end{proof}

A \omterm{def:ml:online-learning-and-drift:rls}{forgetting factor} is therefore an assumption about how fast the truth moves. \Cref{fig:ml:online:forgetting} measures
the trade-off: short memories track regimes but estimate five coefficients from little data; long ones estimate well
the coefficients of regimes that have ended. With regimes of 100 or 500 steps on average, the best memory is about 200
observations ($\lambda = 0.995$), and it recovers 0.75\% and 2.44\% of R-squared out of a possible 4.9\%; with regimes of
2\,000 steps it is 500 observations (3.71\%), and with regimes of 10\,000 steps nothing should be forgotten ($\lambda = 1$,
5.21\% of 5.29\%). Online stochastic gradient descent with a fixed \omterm{def:ml:trees-and-boosting:boosting}{learning rate} is a crude forgetting rule of its own
(0.57\%, 2.05\%, 3.69\% and 4.84\%).

\begin{omfigure}
\begin{tikzpicture}
\begin{semilogxaxis}[width=11cm,height=6cm,xlabel={effective memory $1/(1-\lambda)$ (observations; right: $\lambda = 1$)},
  ylabel={prequential $R^2$ (\%)},tick label style={font=\small},label style={font=\small},xmin=40,xmax=25000,
  ymin=-2,ymax=6,xtick={50,100,200,500,1000,2000,20000},xticklabels={50,100,200,500,1\,000,2\,000,all},grid=major,
  grid style={gray!25},legend columns=4,legend style={font=\footnotesize,at={(0.5,-0.25)},anchor=north,draw=none,
  fill=none,/tikz/every even column/.append style={column sep=6pt}},table/col sep=comma]
\addplot[omThm,very thick,mark=*] table[x=memory,y=D100]{figdata/ml/12-online-learning-and-drift/forgetting.csv};
\addplot[omProp,very thick,mark=square*] table[x=memory,y=D500]{figdata/ml/12-online-learning-and-drift/forgetting.csv};
\addplot[omDef,very thick,mark=triangle*] table[x=memory,y=D2000]{figdata/ml/12-online-learning-and-drift/forgetting.csv};
\addplot[black!70,very thick,mark=diamond*] table[x=memory,y=D10000]{figdata/ml/12-online-learning-and-drift/forgetting.csv};
\legend{$D = 100$,$D = 500$,$D = 2\,000$,$D = 10\,000$}
\end{semilogxaxis}
\end{tikzpicture}

\omcaption{Prequential R-squared of \omterm{def:ml:online-learning-and-drift:rls}{recursive least squares} (after 1\,000 steps of warm-up) against its effective memory, for
streams whose coefficients change every $D$ steps on average. Data: \texttt{ml\_online.forgetting}.}\label{fig:ml:online:forgetting}
\end{omfigure}

\section{Detecting drift}

\begin{definition}[Drift detector, Page--Hinkley test, adaptive windowing]\label{def:ml:online-learning-and-drift:detector}
A \emph{drift detector}\index{drift detector} watches a statistic of a model in production (its errors, its gains, a
feature's distribution) and raises an alarm when the statistic's distribution appears to have changed. The
\emph{Page--Hinkley test}\index{Page--Hinkley test} accumulates the deviations of the monitored value below its running mean,
less a tolerance $\delta$, and alarms when the sum rises more than $h$ above its running minimum (Page, 1954; Hinkley, 1971).
\emph{Adaptive windowing}\index{adaptive windowing} (ADWIN) keeps a window of recent values and drops its older part whenever
some split of the window into two sub-windows has means further apart than a Hoeffding-type bound at confidence $\delta$
(Bifet and Gavaldà, 2007).
\end{definition}

Book~7's CUSUM test (chapter~13) is the third detector. All three watch the same statistic: the deployed model's gain
$g_t = \hat y_t\,y_t$ (its forecast times the outcome), standardised on the warm-up period, whose mean is positive while the
model is right and falls when it is not. Their thresholds are set, on a long stream with no change, to one false alarm a
year (one per 252 steps): $h = 9$ for CUSUM and Page--Hinkley (tolerance 0.1), $\delta = 0.0032$ for ADWIN. Then, on twenty
streams whose coefficients reverse sign at step 2\,000, the detectors raise their first true alarm after 20.8 steps on
average for CUSUM (median 17), 27.4 for Page--Hinkley (median 18.5) and 42.5 for ADWIN (median 26.5); over the 1\,000 steps
before the reversal they raised 74, 78 and 91 false alarms across the twenty streams, close to the 79 their calibration
implies (\cref{fig:ml:online:alarms}).

\begin{omfigure}
\begin{tikzpicture}
\begin{axis}[width=12cm,height=5.6cm,xlabel={step},ylabel={50-step mean of the gain},tick label style={font=\small},
  label style={font=\small},xmin=1600,xmax=2400,ymin=-1.5,ymax=1.9,grid=major,grid style={gray!25},
  xticklabel style={/pgf/number format/1000 sep={}},legend columns=4,legend style={font=\footnotesize,at={(0.5,-0.3)},
  anchor=north,draw=none,fill=none,/tikz/every even column/.append style={column sep=6pt}},table/col sep=comma]
\addplot[gray,thick] table[x=step,y=running]{figdata/ml/12-online-learning-and-drift/gains.csv};
\addplot[only marks,mark=triangle*,mark size=3pt,omThm,restrict expr to domain={\thisrow{detector}}{0:0}]
  table[x=step,y=level]{figdata/ml/12-online-learning-and-drift/alarms.csv};
\addplot[only marks,mark=square*,mark size=2.6pt,omProp,restrict expr to domain={\thisrow{detector}}{1:1}]
  table[x=step,y=level]{figdata/ml/12-online-learning-and-drift/alarms.csv};
\addplot[only marks,mark=*,mark size=2.6pt,omDef,restrict expr to domain={\thisrow{detector}}{2:2}]
  table[x=step,y=level]{figdata/ml/12-online-learning-and-drift/alarms.csv};
\draw[black!70,thick,dashed] (axis cs:2000,-1.5) -- (axis cs:2000,1.9);
\legend{gain,CUSUM alarm,Page--Hinkley alarm,ADWIN alarm}
\end{axis}
\end{tikzpicture}

\omcaption{One stream whose coefficients reverse at step 2\,000 (dashed): the 50-step mean of the deployed model's
standardised gain, and each detector's alarms from 300 steps before the reversal to its first alarm after it (markers at
arbitrary heights, one row per detector). Data: \texttt{ml\_online}.}\label{fig:ml:online:alarms}
\end{omfigure}

\section{Retraining schedules}

\begin{definition}[Retraining schedule]\label{def:ml:online-learning-and-drift:schedule}
A \emph{retraining schedule}\index{retraining schedule} is the rule that decides when a model is refitted and on which
data: on a calendar (every week, on a rolling window), on a detector's alarm, continuously (an online estimator), or by a
combination; it is part of the model and is validated with it.
\end{definition}

\begin{table}[htbp]
\centering\small
\begin{tabular}{@{}lcc@{}}
\toprule
schedule (least squares on the last 500 observations unless stated) & prequential $R^2$ & refits\\
\midrule
never (fitted on the first 1\,000 steps) & $-7.81\%$ & 0\\
every 250 steps & 2.62\% & 75\\
on a Page--Hinkley alarm (refit on the data since the alarm) & $-2.67\%$ & 2\\
\omterm{def:ml:online-learning-and-drift:rls}{recursive least squares}, $\lambda = 0.998$ & 3.24\% & every step\\
\midrule
the truth & 4.62\% & \\
\bottomrule
\end{tabular}
\caption{Four \omterm{def:ml:online-learning-and-drift:schedule}{retraining schedules} on a stream whose coefficients are redrawn every 2\,000 steps on average (19\,000 steps
scored after a warm-up of 1\,000). Data: \texttt{ml\_online.schedules}.}\label{tab:ml:online:schedules}
\end{table}

\Cref{tab:ml:online:schedules} holds a lesson about detectors. The Page--Hinkley detector that caught a sign reversal
within a median 18.5 steps refitted the model twice while the coefficients changed ten times: when coefficients are redrawn rather than reversed, the old
model's gains fall to zero rather than below it, a smaller change than it was tuned to fear, and it waits. The calendar
schedule, dumb and regular, beats it; \omterm{def:ml:online-learning-and-drift:rls}{recursive least squares} with the memory that suits the regimes beats both. A detector
must be tuned to the change that matters, and a schedule that relies on one needs a calendar underneath it.

\begin{method}[Choosing how a model adapts]\label{met:ml:online:method}
\begin{enumerate}
\item Estimate how fast the relation moves: fit the model on rolling windows of several lengths and compare their
prequential scores, or equivalently scan the \omterm{def:ml:online-learning-and-drift:rls}{forgetting factor}.
\item Score every schedule prequentially on history, including the costs of refitting (validation, deployment, turnover).
\item Calibrate detectors on periods without change to a false-alarm rate the desk can act on, and measure their delay on
the changes that matter (a reversal, a fade to zero, a volatility shift).
\item Combine: a calendar refit as the floor, a detector to refit early, and an online estimator where the model is
simple enough to update safely.
\end{enumerate}
\end{method}

\section{Tutorial: the half-life of a model}\label{tut:ml:online-learning-and-drift:tut}

\begin{tutorial}
\textbf{Goal.} Scan the \omterm{def:ml:online-learning-and-drift:rls}{forgetting factor} of \omterm{def:ml:online-learning-and-drift:rls}{recursive least squares} against the regime length, calibrate three detectors
to one false alarm a year and time them on a reversal, and compare \omterm{def:ml:online-learning-and-drift:schedule}{retraining schedules}. \textbf{End state:}
\cref{fig:ml:online:forgetting,fig:ml:online:alarms}, \cref{tab:ml:online:schedules}.
\begin{tutsteps}
\item \textbf{\omterm{def:ml:online-learning-and-drift:rls}{Recursive least squares} with forgetting} (\cref{prop:ml:online:kalman}).
\omcode{code/firm/onlinelearn/firm_onlinelearn.py}{51}{67}{Recursive least squares with a forgetting factor.}\label{lst:ml:online:rls}
\item \textbf{The \omterm{def:ml:online-learning-and-drift:detector}{Page--Hinkley test}.}
\omcode{code/firm/onlinelearn/firm_onlinelearn.py}{104}{123}{The Page--Hinkley test for a fall in the mean.}\label{lst:ml:online:ph}
\item \textbf{A windowed ADWIN.}
\omcode{code/firm/onlinelearn/firm_onlinelearn.py}{126}{155}{Adaptive windowing, checked every few steps.}\label{lst:ml:online:adwin}
\item \textbf{Run} \texttt{ml\_online.forgetting(D)}, \texttt{thresholds()}, \texttt{delays()}, \texttt{schedules()} and
\texttt{fig\_online.py}.
\end{tutsteps}
\textbf{What to change next.} Make the regimes fade their coefficients to zero instead of redrawing them and recalibrate the
detectors for that change; combine the calendar with the detector.
\end{tutorial}

\section{Build: online learning and drift detection}\label{bld:ml:online-learning-and-drift:onlinelearn}

\begin{build}
\bfield{Purpose} Models that adapt at a chosen speed, detectors with a known false-alarm rate, and schedules validated like
the models they refit.
\bfield{Interface} \texttt{drift\_stream(n, p, mean\_regime, r2, seed, flip\_at)}; \texttt{RLS(p, lam, delta)},
\texttt{OnlineSGD(p, lr)}, \texttt{prequential(model, X, y)}; \texttt{CUSUM(mu0, k, h)}, \texttt{PageHinkley(delta, lam)},
\texttt{ADWIN(delta, max\_window, check, min\_side)}, \texttt{alarm\_times}, \texttt{calibrate(make\_for, thresholds,
null\_values, rate)}; \texttt{retrain\_schedule(X, y, policy, every, window, warm, make\_detector, min\_fit)}.
\bfield{Rules} Predict before learning; detectors see only past and present values; thresholds come from calibration on
no-change data, never from the period being judged.
\bfield{Acceptance tests} \texttt{code/firm/onlinelearn/tests/}: RLS with $\lambda = 1$ equals batch least squares; its effective
memory matches $1/(1-\lambda)$ on a planted change; each detector's false-alarm rate on no-change data is within its
calibration and it finds a planted mean shift; the calendar schedule refits at the right steps.
\bfield{Stretch} Forgetting that adapts to the detector's statistic (variable $\lambda$); detectors on feature
distributions (chapter~27).
\end{build}

\begin{omsources}
\item E.~S.~Page, ``Continuous inspection schemes'', \emph{Biometrika} 41(1/2), 1954.
\item A.~Bifet and R.~Gavaldà, ``Learning from time-changing data with adaptive windowing'', \emph{SIAM International
Conference on Data Mining}, 2007.
\item J.~Gama, I.~Žliobaitė, A.~Bifet, M.~Pechenizkiy and others, ``A survey on concept drift adaptation'', \emph{ACM
Computing Surveys} 46(4), 2014.
\item A.~H.~Sayed and T.~Kailath, ``A state-space approach to adaptive RLS filtering'', \emph{IEEE Signal Processing
Magazine} 11(3), 1994.
\item D.~V.~Hinkley, ``Inference about the change-point from cumulative sum tests'', \emph{Biometrika} 58(3), 1971.
\item G.~Widmer and M.~Kubat, ``Learning in the presence of concept drift and hidden contexts'', \emph{Machine Learning} 23,
1996.
\end{omsources}

\section{Exercises}

\begin{exercise}[$\star$]\label{exo:ml:online-learning-and-drift:1}
What are the effective memories of \omterm{def:ml:online-learning-and-drift:rls}{forgetting factors} 0.99, 0.998 and 0.9995? What weight does an observation 500 steps old
get with each, relative to the newest?
\end{exercise}

\begin{exercise}[$\star$]\label{exo:ml:online-learning-and-drift:2}
A detector raises one false alarm per 252 steps. How many false alarms should 20 streams of 1\,000 steps raise? Compare with
the chapter's counts.
\end{exercise}

\begin{exercise}[$\star$]\label{exo:ml:online-learning-and-drift:3}
A CUSUM with reference 0 and allowance $k = 0.1$ watches standardised gains whose mean falls from 0 to $-0.38$ at the
reversal (the chapter's streams). At what average rate does its statistic grow after the change, how many steps does it
need from zero to pass $h = 9$, and why is the measured median delay shorter?
\end{exercise}

\begin{exercise}[$\star\star$]\label{exo:ml:online-learning-and-drift:4}
Why does the best memory grow with the regime length, and why is it not proportional to it?
\end{exercise}

\begin{exercise}[$\star\star$]\label{exo:ml:online-learning-and-drift:5}
Why did the Page--Hinkley schedule refit only twice, and what would you change?
\end{exercise}

\begin{exercise}[$\star\star$]\label{exo:ml:online-learning-and-drift:6}
\emph{Find the flaw.} ``We chose the detector's threshold so that it caught last year's two regime changes within a week,
and it would have.''
\end{exercise}

\begin{exercise}[$\star\star\star$]\label{exo:ml:online-learning-and-drift:7}
\emph{Coding.} Run \texttt{ml\_online.schedules} with regimes of 500 steps on average (\texttt{schedules(D=500)}). Report the
four schedules and explain what changed.
\end{exercise}

\begin{exercise}[$\star\star\star$]\label{exo:ml:online-learning-and-drift:8}
Show that RLS with $\lambda = 1$ and $P_0 = \delta I$ gives the ridge estimator with penalty $1/\delta$, and hence ordinary
least squares as $\delta\to\infty$.
\end{exercise}

\section{Problem: The Half-Life of a Model}

\begin{problem}[{Weekend problem --- how fast to forget}]\label{pb:ml:online-learning-and-drift:1}
The chapter's stream: five features, a signal worth 5\% of the variance, coefficients redrawn every $D$ steps on average.

\medskip\noindent\textbf{Part I --- Forgetting.}
\begin{enumerate}
\item What is the truth's R-squared, and why is no model close to it for short regimes?
\item Which \omterm{def:ml:online-learning-and-drift:rls}{forgetting factor} is best for $D = 100$, 500, 2\,000 and 10\,000, and what does it score?
\item What does online SGD score, and why does it behave like forgetting?
\item What does \cref{prop:ml:online:kalman} say a \omterm{def:ml:online-learning-and-drift:rls}{forgetting factor} assumes?
\end{enumerate}

\medskip\noindent\textbf{Part II --- Detectors.}
\begin{enumerate}[resume]
\item What statistic do the detectors watch, and why the gain rather than the squared error?
\item What thresholds give one false alarm a year?
\item What are their mean and median delays after a reversal?
\item How many false alarms did they raise before the reversal, and is that too many?
\end{enumerate}

\medskip\noindent\textbf{Part III --- Schedules.}
\begin{enumerate}[resume]
\item What do the four schedules score with regimes of 2\,000 steps?
\item Why does the never-refitted model score below zero?
\item Why did the detector-driven schedule fail?
\item Why does \omterm{def:ml:online-learning-and-drift:rls}{recursive least squares} win?
\end{enumerate}

\medskip\noindent\textbf{Part IV --- The verdict.}
\begin{enumerate}[resume]
\item State the \emph{named result}: the \omterm{def:ml:online-learning-and-drift:rls}{forgetting factor} that minimises the out-of-sample loss as a function of the regime
length, and each detector's delay at one false alarm a year.
\item What schedule would you run for a short-horizon model on a changing venue?
\item How would you estimate $D$ on real data?
\item What does a false alarm cost, and what does a missed change cost?
\item How does this chapter connect to Book~7's signal half-life?
\item What should a monitoring system log for each alarm?
\item How would \omterm{def:ml:online-learning-and-drift:drift}{covariate shift} show up in these statistics?
\item In one sentence: how fast should a model forget?
\end{enumerate}
\end{problem}

\section{Interview questions}

\begin{interviewq}[$\star$ \iqroles{mle, researcher}]\label{iq:ml:online-learning-and-drift:1}
What is the difference between \omterm{def:ml:online-learning-and-drift:drift}{concept drift} and \omterm{def:ml:online-learning-and-drift:drift}{covariate shift}? Give a market example of each.
\end{interviewq}

\begin{interviewq}[$\star\star$ \iqroles{researcher}]\label{iq:ml:online-learning-and-drift:2}
Derive \omterm{def:ml:online-learning-and-drift:rls}{recursive least squares} with a \omterm{def:ml:online-learning-and-drift:rls}{forgetting factor}, and say what the \omterm{def:ml:online-learning-and-drift:rls}{forgetting factor} means.
\end{interviewq}

\begin{interviewq}[$\star\star$ \iqroles{mle}]\label{iq:ml:online-learning-and-drift:3}
How do you set the threshold of a \omterm{def:ml:online-learning-and-drift:detector}{drift detector}, and how do you know it works?
\end{interviewq}

\begin{interviewq}[$\star\star$ \iqroles{researcher, trader}]\label{iq:ml:online-learning-and-drift:4}
Your model is retrained monthly. How would you decide whether to retrain weekly, daily, or continuously?
\end{interviewq}

\begin{interviewq}[$\star\star$ \iqroles{mle}]\label{iq:ml:online-learning-and-drift:5}
Explain the \omterm{def:ml:online-learning-and-drift:detector}{Page--Hinkley test} and one weakness of it.
\end{interviewq}

\begin{interviewq}[$\star\star\star$ \iqroles{researcher}]\label{iq:ml:online-learning-and-drift:6}
Show that exponential forgetting is a Kalman filter for a random-walk coefficient, and what that implies for choosing
$\lambda$.
\end{interviewq}
